\documentclass[journal]{IEEEtran}

\usepackage[numbers,sort&compress]{cite}
\usepackage{graphicx}
\usepackage{amsmath}
\usepackage{amssymb}
\usepackage{url}
\usepackage[colorlinks=true,linkcolor=black,citecolor=black,urlcolor=black]{hyperref}
\usepackage{booktabs}
\usepackage{caption}
\usepackage{subcaption}
\usepackage{array}
\usepackage{algorithmic}

\graphicspath{{./figures/}}

\ifCLASSINFOpdf
  \DeclareGraphicsExtensions{.pdf,.png,.jpg}
\else
  \DeclareGraphicsExtensions{.eps}
\fi

\begin{document}


\title{RDAF-Net: A Single-Stage Rotated Dual Alignment Feature Network for Remote Sensing Object Detection}

\author{\IEEEauthorblockN{Author Names}}
\IEEEauthorblockA{Affiliation, City, Country}

\markboth{IEEE Transactions on Geoscience and Remote Sensing}%{
\MakeLowercase{\textit{RDAF-Net: A Single-Stage Rotated Dual Alignment Fe...}}}

\maketitle

\begin{abstract}
The field of remote sensing object detection has experienced rapid development in recent years, demonstrating significant value in various applications. However, with the increasing demand for detection and the complexity of scenarios, oriented object detection has become a major challenge: firstly, the rotation-sensitive feature representation of traditional convolutional networks is inaccurate, making it difficult to accurately model the slender shape features of high aspect ratio targets; secondly, the mismatch between the static label allocation strategy and the feature space of dynamic rotating targets leads to bias in training sample selection under extreme aspect ratio scenarios. To address these issues, this paper proposes a single-stage Rotated Dual Alignment Feature Network (RDAF-Net) that collaboratively enhances detection accuracy through feature representation optimization and dynamic matching mechanisms. Specifically: 1) a Dynamic Refined Rotation Convolution Module (DRRCM) is designed, which uses a Data-Enhanced Spatial Attention Module (DESAM) to predict dynamic rotation convolution kernel parameters, achieving adaptive feature alignment according to the target's rotation angle; 2) an Anchor Refinement Feature Alignment Module (ARFAM) is constructed, which solves the feature-space misalignment problem by adjusting the feature sampling position guided by preset anchor boxes through regression correction; 3) a Shape-Aware Quality Assessment (SAQA) strategy is proposed, which dynamically adjusts the IoU threshold based on target shape information (SI) and optimizes sample matching quality in conjunction with the distance-weighted shape center. Experimental results on three commonly used remote sensing target datasets (i.e., HRSC2016, DOTA, and UCAS-AOD) demonstrate the effectiveness of our proposed framework.
\end{abstract}

\IEEEpeerreviewmaketitle


\section{ Introduction }

In recent years, with the rapid development of deep learning, the ability of convolutional neural networks to extract feature information has further promoted the development of object detection in the field of remote sensing [1-5]. Remote sensing object detection, as one of the core research directions in the field of computer vision, has important application value. Its main task is to accurately identify targets of interest in given remote sensing images and classify and locate them [6-8]. In remote sensing images, targets are mostly oriented differently and arranged closely [9, 10], which makes the adaptability of general object detection frameworks poor. The horizontal anchor boxes often contain a lot of background interference and the problem of multiple targets being selected simultaneously, leading to a significant decline in detection performance. Therefore, how to effectively detect oriented objects in remote sensing images has become a current research hotspot. Recently, some studies have focused on improving the representation of oriented bounding boxes in remote sensing object detection. This is mainly achieved by developing specialized detection frameworks, such as R3Det[11], Rotated RetinaNet[12], and RoI Transformer[13], as well as oriented box encoding techniques, such as sliding vertex offset[14], short side offset[15], and midpoint offset box encoding[16]. In addition, various loss functions, including CSL[17], KLD[18], and KFIoU[19], have been proposed to further improve the performance of these methods.

Although the above methods exhibit high precision in detecting most targets, the detection accuracy for certain specific types of targets (e.g., bridges, ports, ships, etc.) is still unsatisfactory. We believe that the root of the problem lies in the fact that existing methods overlook the differences in shape information between different types of targets, especially the distinction between high aspect ratio targets and regularly shaped targets. Here, the shape information we mentioned refers to the ratio of the long side to the short side of the ground-truth bounding box of the target object, which can be intuitively represented by a mathematical formula:, where \textit{w} and \textit{h} represent the width and height of the ground-truth bounding box, respectively, and the value range of \textit{SI} is [0,1]. We classify all objects into two categories, with1＞\textit{SI}＞for regularly shaped targets and\textit{SI}≤for high aspect ratio targets, as shown in Figure 1.

(a) Regularlyshaped target(b) High aspect ratio target

Fig.1. Shape information distinction. The green box represents the ground-truth bounding box. \textit{w}, \textit{h}denote the width and height of the ground-truth bounding box, SI represents the shape information, \textit{SI} greater than one-half indicates a regularly shaped target, and\textit{ SI} less than one-half is expressed as a high aspect ratio target.

Existing methods perform well in detecting regularly shaped targets, but when faced with rotated targets with large aspect ratio variations, the detection accuracy significantly decreases. There are two main reasons for this:

(1) Feature misalignment: The convolutional features of traditional backbone networks are usually aligned based on fixed receptive field orientations, making it difficult to adapt to directionally oriented targets with high aspect ratio differences, resulting in poor feature extraction effects. Even if convolutional alignment or anchor box alignment operations are introduced in subsequent steps, they cannot compensate for the loss of local edge information of the target caused by the initial fixed convolution method, thereby affecting the overall quality of feature extraction. This is because all subsequent operations, such as feature fusion and resampling, are based on the feature maps extracted by the initial backbone network. Therefore, in the object detection framework, the backbone network that initially extracts target features is crucial for improving model accuracy.

(2) Sample imbalance: In the process of anchor box regression, high aspect ratio rotated targets are extremely sensitive to angle regression. Even a slight angular deviation can cause a significant increase in the deviation between the predicted and ground truth boxes, especially when the shape information value is small, as shown in Figure 2. This situation leads to an increase in false negatives during sample selection. Even with high classification scores, due to poor regression performance, Intersection over Union (IoU)between the predicted and ground truth boxes is below the preset threshold, causing targets that should be positive samples to be misclassified as negative samples. This misclassification results in an imbalance between positive and negative samples, which in turn negatively affects the overall detection performance.

(c)                            (d)

Fig. 2. In (c) and (d), the green boxes represent the ground-truthbounding boxes, while the burgundy and yellow boxes represent the predicted boxes rotated by 15° and 30°, respectively. The overlapping areas of the burgundy and yellow boxes indicate IoU with the ground-truth boundingboxes. It can be clearly observed that the impact of the rotation angle on ships with a high aspect ratio is much greater than that on planes with regular shapes.

To address this challenge, we approach from two aspects: high-quality feature extraction and rational label allocation. We believe that high-quality feature extraction is the foundation of oriented object detection, while rational label allocation can further enhance the detection performance. We propose a single-stage Rotated Dual Alignment Feature Network (RDAF-Net), which consists of two modules and one matching strategy: Dynamic Refined Rotation Convolution Module (DRRCM), Anchor Refinement Feature Alignment Module (ARFAM) and Shape-Aware Quality Assessment (SAQA) matching strategy. Specifically, the DRRCM in the backbone network can accurately predict the weight and angle of the rotational convolution kernel using the Data-Enhanced Spatial Attention Module (DESAM). Subsequently, the predicted parameters are combined to generate a convolution kernel that adaptively adjusts according to the pose information of the oriented target, achieving accurate alignment with the target features and producing direction-sensitive feature maps. The ARFAM in the detection head quickly generates high-quality modified prediction anchor boxes on the direction-sensitive feature map through the regression branch, serving as guidance to dynamically adjust the position of the feature sampling points, thus further achieving precise feature alignment. High-quality rotation-invariant features are extracted through the collaborative work of the double alignment convolution module. In the label allocation process, based on the high-quality feature maps and modified prediction anchor boxes generated above, the SAQA method is used to dynamically adjust the IoU threshold according to the shape information of the target for training sample selection, and the quality information is added to the selected positive samples using the shape center distance weighting method to optimize training sample selection. Extensive experiments on commonly used datasets (such as HRSC2016, DOTA, and UCAS-AOD) show that our method can maintain a good detection performance.

The main contributions of this paper are summarized as follows:

1) We propose a novel object detection framework, RDAF-Net, to generate high-quality rotation-invariant features and regression prediction anchor boxes and couple efficient sample selection strategies to obtain excellent detection performance.

2) A flexible dynamic rotation convolution module is proposed, which can be easily embedded into the backbone networks of many detectors to extract high-quality basic features for oriented targets.

3) A reasonable sample matching strategy, using sample shape information and potential sample quality to optimize training samples, thus solving the inconsistency between classification and regression.

\section{ Related Work }

\subsection{ Remote Sensing Object Detection }

Objects in natural scenes mostly exhibit characteristics such as arbitrary orientations, significant scale variations, and dense distributions [10,20]. General object detection methods using horizontal anchors face the same issue, as their horizontal bounding box localization fails to accurately describe the pose diversity of remote sensing targets. To enable detectors to precisely detect oriented objects, researchers have proposed methods like angle prediction and preset rotated anchor boxes to improve regression, albeit with increased computational complexity.Ma et al. [21] proposed the Rotated Region Proposal Network (RRPN), which introduces rotated candidate boxes into the Region Proposal Network (RPN) architecture to detect tilted or non-axis-aligned targets. They adopt \textit{(x,y,w,h,θ)} as positional parameters for rotated bounding boxes, generating anchor boxes with varying scales, aspect ratios, and rotation angles. Fine-tuning these parameters yields angle-aware candidate regions.Ding et al. [13] designed the RoI Transformer, which uses a lightweight fully connected layer to learn geometric parameters of Rotated Regions of Interest (RRoI) from Horizontal Regions of Interest (HRoI), effectively avoiding the need to design numerous RRoIs for oriented object detection. Xu et al. [14] generated rotated anchor boxes by sliding the four vertices of horizontal anchor boxes. Xie et al. [16] proposed the midpoint offset method, displacing the midpoints of the upper and right edges of horizontal anchor boxes and refining them to obtain more accurate rotated anchor boxes.Compared to these methods, the proposed RDAF-Net improves the feature representation of the entire network. First, the DRRCM in the backbone network generates direction-sensitive features, significantly reducing the complexity of modeling oriented targets. Second, the ARFAM in the detection head further refines features to efficiently extract discriminative and semantic features of target boundaries. Finally, the high-quality feature maps are dynamically assigned labels through the SAQA matching strategy, effectively alleviating the inconsistency between regression and classification.

\subsection{ Feature Extraction for Object Detection }

Feature extraction is a critical component in object detection, as the quality of features extracted by the network directly determines the overall performance of the model. Azimi et al. [22] designed an image cascade network that generates four input images of different scales through bilinear interpolation, enabling the model to extract multi-scale features. Liu et al. [23] added convolutional modules after each stage of Swin-Transformer [24], combining the local properties of CNNs with the global properties of Transformers to extract richer foundational features. Sun et al. [25] introduced a feedback connection into the feature fusion network to enhance the model's ability to represent multi-scale features. To extract refined features, Zheng et al. [26] proposed a Feature Pyramid Transformer (FPT), which transforms, rearranges, concatenates, and convolves features within the feature pyramid. However, these methods still struggle to achieve satisfactory detection performance when handling complex and diverse oriented targets.Unlike the aforementioned approaches, we propose a dual-layer feature alignment technique. First, the DRRCM adaptively adjusts convolutional kernels based on the pose information of oriented targets to achieve preliminary alignment of target features. Second, the ARFAM dynamically adjusts the positions of feature sampling points guided by refined predicted anchor boxes, further enhancing convolutional alignment. These two convolutional strategies work synergistically at different stages, effectively avoiding the misalignment between axis-aligned convolutional features and arbitrarily oriented targets, thereby extracting high-quality rotation-equivariant features.

\subsection{ Sample Selection for Object Detection }

In detection networks, high-quality feature extraction is undoubtedly crucial as it provides a solid foundation for subsequent model training. However, the sample selection process is equally critical, exerting a decisive influence on both training efficacy and final detection performance. Most object detectors conventionally adopt fixed IoU threshold-based matching strategies to select positive and negative samples during training [27]. However, this fixed metric approach often demonstrates poor adaptability when confronted with significant inter-sample variations—such as target scale diversity, shape variability, and background complexity—leading to suboptimal detection outcomes [28]. Consequently, dynamic sample selection strategies have gained increasing attention. Li et al. [29] proposed an Adaptive Points Assessment and Assignment (APAA) scheme, which dynamically selects the top-k samples with the highest quality scores as positive training samples by evaluating point sets across four dimensions: classification quality, localization accuracy, orientation consistency, and point correlation. Ming et al. [30] introduced Dynamic Anchor Learning (DAL), which comprehensively evaluates the localization potential of anchors through dynamic assignments based on predefined matching scores. Zhang et al. [31] developed a Learning-to-Match (LTM) method under the principle of maximum likelihood estimation, enabling targets to dynamically select optimal anchors. This paper employs the SAQA method for label assignment, which comprehensively integrates the target's shape information and center-point information to enable more flexible sample selection.

\section{ Methodology }

The overall architecture of RDAF-Net is illustrated in Figure 3. First, it adopts the DRRCM as the backbone network to dynamically adjust convolutional kernels for aligning arbitrarily oriented objects, thereby extracting preliminary aligned features. Subsequently, the ARFAM performs anchor refinement to guide convolutions in acquiring more precise aligned features. Finally, leveraging the high-quality feature maps and refined anchor boxes generated through the above stages, the SAQA strategy dynamically assigns labels and evaluates sample quality, ensuring optimal matching between candidate boxes and their corresponding ground-truth labels. This approach effectively avoids the inconsistency between regression and classification, thereby enhancing detection performance. The detailed implementation of RDAF-Net is elaborated below.

Fig. 3. Proposed RDAF-Net architecture.

\subsection{ Dynamic Refined Rotation Convolution Module }

In most existing remote sensing object detectors, the convolutional structures used in the backbone networkadopt axis-aligned or preset fixed rotation angles for feature extraction of targets. However, objects in natural scenes are often placed at arbitrary angles. Therefore, standard convolutional kernels struggle to precisely match the contours of non-axis-aligned targets, making it difficult to effectively extract high-quality features from these arbitrarily oriented objects.

In order to avoid the fixed convolution mode of standard convolution and enhance the representation ability of the target, thus achieving accurate object detection. We propose a Dynamic Refined Rotation Convolution Module (DRRCM). The overall structure of DRRCM is shown in Figure 4 below.

Fig.4. Illustration for the DRRCM module.

We designed a Data-Enhanced Spatial Attention Module (DESAM) to generate a spatial mask through pooling, concatenation, convolution transformation and sigmoid activation function to weight the fused features and highlight the important spatial regions. Then, the weighted fused features are average pooled and input into the kernel angle prediction branch and the kernel weight prediction branch. This module can make the network more accurately focus on the key feature positions in rotating object detection, and then accurately generate the predicted weights and angles of the rotating kernels.

Specifically, we first efficiently extract spatial relationships by applying channel-wise average pooling and max pooling (denoted asand, respectively) to the feature map \textit{F }obtained from depthwise convolution:

	(1)

Where \textit{S}\textit{avg} and \textit{S}\textit{max} are the spatial feature descriptors obtained through average pooling and max pooling, respectively. To enable information interaction between different spatial descriptors, we concatenate the spatially pooled features and employ a convolutional layer to transform the pooled features (with 2 channels) into \textit{C}\textit{in} spatial attention maps.

	(2)

Where \textit{C}\textit{in} represents the number of input channels. For each spatial attention map, a sigmoid activation function is applied to obtain an individual spatial maskfor each convolutional kernel.

	(3)

Then, the spatial masks are used to weight the features after depthwise convolution, which are subsequently compressed into a C-dimensional feature vector through global average pooling.

	(4)

Then, the pooled feature vectors are sent into two branches respectively.

The first branch predicts the rotation kernel angle, where the feature vector is input into this branch and processed through \textit{Dropout}, a linear layer, \textit{Softsign} activation, and multiplication by a scale factor to obtain a set of angles :

	(5)

Where is the linear layer without a bias term, ensuring that angle prediction depends solely on variations in the input features and avoids learning biased angles.  is a scaling factor used to expand the rotation range, and the parameter  is used to adjust the angle range, with a default value of 40.

	The second branch is for rotating kernel weight prediction. By inputting the feature vectors into this branch and through \textit{Dropout}, a linear layer, and \textit{Sigmoid} activation, a set of weightsare obtained.

	(6)

	Where  is a bias set for the linear layer to improve the flexibility of the model.

DESAM is initialized from a truncated normal distribution with a mean of zero and a standard deviation of 0.2, in order to assist the model in converging more rapidly and to reduce instability at the beginning of training. The implementation of the rotational convolution function is elaborated on below.

The rotation angleparameter generated by the above DESAMprediction reparameterizes the weights inside the convolution kernel, allowing the convolution kernel to dynamically adjust according to different input feature maps, achieving adaptive rotation:

	(7)

	(8)

	Whichrepresents the coordinates of the original sampling points.represents the new sampling point coordinates after the original sampling pointis rotated counterclockwise by an angleto achieve the convolution kernel and feature alignment.represents the reparameterized convolution kernel.represents the bilinear interpolation, which is used to calculate the weight value of the new position after the rotation of the  convolution kernel.

The reparameterized convolution kernel is then multiplied by the corresponding \textit{λ}\textit{i}weight and summed, and then convolved with the input feature map to finally generate high-quality direction-aware features \textit{Y}:

	(9)

Through the DRRCM, the convolution kernel can adaptively adjust according to the different orientations of objects in the input feature map, thereby efficiently capturing the features of multi-directional objects in the image. Especially when detecting objects in aerial images that are densely arranged and have large scale differences, high-quality feature extraction is crucial for accurate classification and precise localization.

\subsection{ Anchor Refinement Feature Alignment Module }

Based on the high-quality feature maps generated by DRRCM in the backbone network, Anchor Refinement Feature Alignment Module(ARFAM) is employed to further refine anchor points through a regression operation. The refined anchor parameters are then used to compute an offset field, enabling the dynamic adjustment of aligned convolution sampling points. This process generates feature representations that are more precisely aligned with the target object, as shown in Figure 5.

Fig.5. Illustration of the ARFAM module.

First, unlike most dense anchor sampling methods, we only preset a single initial square anchor at each position on the feature map. This anchor is then refined into a high-quality oriented anchor through the regression branch, thus reducing the need to preset a large number of anchors on the feature map, which helps to reduce computational complexity. The offset for predicting the anchor box regression target is as follows:

 ; 

 ; 	(10)

	Where (\textit{x},\textit{y},\textit{w},\textit{h},\textit{θ}) represent the center coordinates, width, height, and angle parameters of the initial anchor. (\textit{x}\textit{g},\textit{y}\textit{g},\textit{w}\textit{g},\textit{h}\textit{g},\textit{θ}\textit{g}) represent the parameters (same as above) of the ground-truth bounding box. (\textit{Δx}\textit{g},\textit{Δy}\textit{g},\textit{Δw}\textit{g},\textit{Δh}\textit{g},\textit{Δθ}\textit{g}) represent the offsets between the ground-truth bounding box and the initial anchor. By regressing these offsets, the model can adjust the initial anchor to a corrected predicted anchor box that is closer to the ground-truth bounding box.\textit{R(θ)} represents the rotation transformation matrix used to convert the center coordinates of the ground-truth bounding box to the coordinate system relative to the initial anchor.\textit{K}is the scaling factor used to adjust the angle value, ensuring that the rotation angle remains within a reasonable range.

	Secondly, in order to achieve feature extraction for the oriented target, we adjust the feature sampling point position to realize dynamic convolution alignment guided by the modified prediction anchor box. This method adds an offsetcalculated from the modified prediction anchor box to the original sampling points of the standard convolution:

	(11)

Here,represents the position of the sampling point for predicting the anchor box. denotes the position of the conventional sampling point for standard convolution, with andrepresenting the two-dimensional coordinates and relative offset of this sampling point, respectively. \textit{R} signifies the regular gridof standard convolution; for instance,would indicate a 3×3 convolution kernel with a dilation rate of 1.

	Then, the dynamic alignment convolution combines the offset and input features x such that the sampled point location can be adjusted based on the shape and orientation of the predicted anchor box to better match the geometry of the actual target.

	(12)

Whererepresents the value of the output feature map at position \textit{p}.is the weight of the convolutional kernel at position.represents the input feature map.

The model is able to extract high-quality feature representations from the input image through a dual-layer feature alignment technique. Its backbone network dynamically refines the preliminary feature extraction of rotational convolution, and the anchor point refinement feature alignment convolution in the detection head further refines the features.These features not only contain the basic information of the target but are also enhanced at the detail level, providing an effective feature basis for subsequent positioning.

\subsection{ Shape-Aware Quality Assessment }

In previous sections, we introduce the extraction of sensitive features with coding direction information in RDAF-Net and its further alignment feature refinement process after anchor box correction, but in the actual training process of the model, it is found that there are still problems of inconsistency between regression and classification tasks, that is, high classification scores cannot guarantee the accurate positioning of detection. This problem has been widely studied in many articles [32–34], and some discussions are traced back to the uncertainty of bounding box regression and localization [34]. We believe the bias between classification and regression mainly comes from unreasonable sample selection methods, and we further solve this problem from the perspective of utilizing target shape information and evaluating sample quality.

Most of the existing detectors usually select positive anchors for training according to a fixed IoU threshold between the anchor and the ground-truth bounding box [27]. However, such sample selection method often ignores the shape information of the target and fails to make potential distinction in quality of the selected positive samples.

To address the above issue, we employ the shape-aware quality assessment (SAQA) method in the training stage. The implementation of this method is introduced in detail in the following parts of this section.

Specifically, firstly, the IoU threshold is adaptively adjusted according to the shape information of the target, and dynamic sample selection is realized. The formula is as follows:

(13)

Whererepresents the IoU threshold for dynamically selecting samples.represents the mean IoU value of the candidate samples. represents the standard deviation of the IoU values of the candidate samples.  is the weighting parameter used to control the influence of aspect ratio on the weight factor. \textit{SI} represents the shape information of the target.

Then, the selected positive samples are refined by introducing the shape-center distance weighting to evaluate the location and add the quality information. In detail, the shape-center distance weighting value  is calculated using Euclidean distance from sample point to object center and the shape information of the object.

(14)

Where (\textit{x}\textit{i}\textit{,y}\textit{i}) represents the center coordinates of the ground-truth bounding box, (\textit{x}\textit{j}​,\textit{y}\textit{j}​) represents the coordinates of the sample point, and \textit{h}\textit{i}​ and \textit{w}\textit{i}​ represent the height and width of the ground-truth bounding box, respectively.

Next, after obtaining the distance-weighted value of the shape center, the quality score \textit{Q}\textit{ij }of the positive sample is calculated.

(15)

In this way, the quality of samples is distinguished by introducing shape-centered distance weighting. The smaller the distance value, the closer the sample point is to the target center, and the higher its quality score. Therefore, using this distance value allows for a more accurate evaluation of each positive sample's quality, thereby optimizing the sample selection process.

The SAQA method fully utilizes the geometric properties of the target and the potential quality of the samples, ensuring the reasonable selection of positive samples during the training process. This improves the consistency between regression and classification, thereby enhancing the overall performance of the detector.

\section{ Experiments and Analysis }

\subsection{ Datasets }

The experiments were conducted on three publicly available high-resolution remote sensing datasets: HRSC2016, DOTA, and UCAS-AOD. The ground-truth bounding boxes in these datasets are used for oriented bounding box annotations.

The HRSC2016 dataset, released by Liu et al. [35] in 2016, is a dataset for optical remote sensing image ship detection. It consists of images collected from Google Earth of six iconic harbor locations, containing only two types of scenes: offshore and nearshore ships. The dataset includes 1,061 images and 2,976 object instances, with image spatial resolutions ranging from 0.4 to 2 meters, and pixel sizes ranging from 300×300 pixels to 1500×900 pixels. In our experiments, all images were resized to 512×800. The training set consists of 436 images, and the validation set contains 181 images, both used for training, while the remaining 444 images are used for testing.

The DOTA dataset, proposed by Xia et al. [36] in 2018, is a large-scale benchmark dataset for aerial image object detection tasks. The dataset includes 2,806 high-resolution aerial images with original image resolutions ranging from 800×800 to 4000×4000 pixels, containing 188,282 finely annotated instances. It covers 15 object categories: airplane (PL), baseball field (BD), bridge (BR), ground track field (GTF), small vehicle (SV), large vehicle (LV), ship (SH), tennis court (TC), basketball court (BC), tank (ST), soccer field (SBF), roundabout (RA), harbor (HA), swimming pool (SP), and helicopter (HC). We cropped the original images into 1024×1024 pixel tiles with a step size of 824 pixels for single-scale model training. To improve the model's multi-scale detection capability, we performed scale transformations on the original images at three scaling ratios: 0.5, 1.0, and 1.5. Then, the same step size strategy of 512 pixels was applied at each scale to generate standardized 1024×1024 pixel tiles.

The UCAS-AOD dataset, released by Zhu et al. [37] in 2015, is an aerial dataset for oriented aircraft and vehicle detection, containing 1,510 images collected from Google Earth. It includes a total of 1,000 images of airplanes and 510 images of vehicles, with an approximate resolution of 659×1280 pixels and 14,596 instances. The dataset has not been officially divided by the authors. We randomly split it into a training set of 755 images, a validation set of 302 images, and a test set of 453 images.

\subsection{ Implementation Details }

In the proposed RDAF-Net, the backbone network adopts ResNet-50 [38], and we introduce our designed DRRCM to extract features within its structure. We use feature pyramids from P3 to P7 for detecting multi-scale targets. For each feature point on the feature map, only one regression anchor is set, and the scale is four times the total stride size. Data augmentation includes random flipping and random rotation. Our weight decay and momentum are set to 0.0001 and 0.9, respectively. The model is trained using the SGD optimizer with an initial learning rate of 0.0025. The model is trained for 36 epochs on the HRSC2016 dataset, 12 epochs on the DOTA dataset, and 72 epochs on the UCAS-AOD dataset. Ablation studies are conducted on the HRSC2016 dataset, as it mainly focuses on detecting ships in remote sensing images, which have significant aspect ratio and scale variations, making it more suitable for testing the detection performance of our model. In the ablation study, all images are resized to 512×800 for training. To ensure a fair comparison with other methods, the mAP metric defined in the PASCAL VOC 2007 challenge [39] is used across all three datasets. We use the MMDetection [40] toolbox to train the model on an RTX 2080Ti GPU, with a total batch size of 2. Each experiment is conducted multiple times, and the stable value is taken as the final result.

\subsection{ Ablation Studies }

\textbf{1) Evaluation of different components:}

Table 1: Influence of each component of RDAF-Net.

To analyze the impact of different components designed in RDAF-Net, we conducted experiments on the HRSC2016 dataset using a controlled variable approach. All experiments were performed with the same settings to ensure the rigor of the experiments. The experimental results for each component are shown in Table 1. Using only the DRRCM, the detection performance improved by 1.79\% compared to the baseline, indicating that DRRCM can more effectively extract high-quality feature representations for accurate detection performance. When both DRRCM and ARFAM were used, the detection performance reached 84.89\% mAP, an improvement of 9.67\%, which demonstrates that even with just one preset anchor point, guiding the convolutional feature alignment using high-quality corrected anchors can further efficiently extract high-quality feature representations while optimizing the anchor points. This indicates that these two methods not only do not conflict but also progressively extract higher-quality feature representations in a more effective manner, significantly improving detection performance. When all three components—DRRCM, ARFAM, and SAQA—were used, the detection performance improved by an additional 4.69\%. It is clear that in the label assignment process, SAQA made reasonable sample selections based on the high-quality feature map and optimized anchor boxes, ultimately achieving 89.58\% mAP, an improvement of 17.94\% over the baseline, proving the effectiveness of our framework.

\textbf{2) Effectiveness evaluation of DRRCM:}

Table 2: Ablation study on the proposed DRRCM.

Note: ("R-50-DRRCM\_S4" indicates that we replaced all 3×3 standard convolution modules in Stage 4 of ResNet-50 with DRRCM. The "+" symbol denotes the addition of the required replaced "Stage.")

To evaluate the effectiveness of our designed DRRCM, we conducted comparative experiments on the HRSC2016 dataset for the backbone integrated with different DRRCMs. In the experiment, we compared the performance of the standard convolution module based on ResNet-50 backbone with several different versions of the backbone integrated with DRRCM. The experimental results are shown in Table 2. We replaced all 3 × 3 standard convolution modules in 3 residual blocks of Stage4 of ResNet-50 with DRRCM, and it can be seen that the detection effect was improved by 0.29\%, which indicates that the DRRCM we proposed can better extract object features. Further, by replacing all 3 × 3 standard convolution modules in 3 and 6 residual blocks corresponding to Stage4 and Stage3 of ResNet-50 with DRRCM, the detection effect was further improved to 72.66\%. This fully demonstrates the continuous improvement of detection performance brought by the multi-level optimization capture of target features with the expansion of DRRCM integration range. Finally, when we replaced all 3 × 3 standard convolution modules in 3, 6 and 4 residual blocks corresponding to the last three stages of ResNet-50 with DRRCM, the final detection effect reached 73.43\% mAP, which is a 1.79\% improvement over the baseline backbone without replacement. These experimental results fully verify the advantages of DRRCM in improving object detection performance.

\textbf{3) Evaluation of the effect of ARFAM:}

Table 3: Ablation study for ARFAM.

We explored the effects of using different numbers of ARFAM on the network design. The experimental results are shown in Table 3. Through experiments, we conclude that ARFAM can improve the detection performance of RDAF-Net, but more modules are not better. The number of ARFAMs is set to 2, and the detection effect is the best, reaching an mAP of 89.58\%. When the number of ARFAMs continues to increase from 2 to 3–4, the detection effects decrease by 0.15\% and 0.64\%, respectively, and it can be seen that the performance of the model shows a trend of slow decline. We speculate that a deeper network structure does not bring the expected improvement for small object detection that requires a small receptive field. On the contrary, a deeper network may lead to an excessively large receptive field, so that the spatial information of small targets is over-fused, and important local details are lost, thereby affecting the detection effect of small targets.

\section{ Results and Analysis }

Table 4: Comparison experiments of HRSC2016.

Note: ('aug' represents data augmentation; 'Anchor' refers to the number of preset anchor points at each position on the feature map.)

1) Results on HRSC2016: The HRSC2016 dataset contains a large number of rotated ship images with high aspect ratio, multi-scale, and arbitrary orientation, which can fully verify the detection performance of our model for high aspect ratio-oriented targets. Our method achieves competitive performance on the HRSC2016 dataset. As shown in Table 4, using R-101-DRRCM as the backbone network and adjusting the input image to 512 × 800 pixels, our method achieves the highest mean average precision (mAP) of 90.05\%. Even when using the lighter R-50-DRRCM, our method still achieves an mAP of 89.58\%. It is worth noting that our method uses only one square anchor at each position of the feature map, but it still outperforms frameworks that preset a large number of rotated anchors at each point of the feature map. For example, R2CNN presets 21 anchors and R3Det presets 126 anchors. Compared with their best detection results, we use only one anchor and do not use data augmentation strategies for training and testing, which achieves an increase of 16.51\% and 0.32\%, respectively. These results show that it is not necessary to preset a large number of rotated anchor boxes of different scales for oriented object detection. More importantly, it is essential to extract high-quality basic features and optimize high-quality prediction anchor boxes, and on this basis, select reasonable sample training for target recognition. Some qualitative results are shown in Fig. 6.

Figure 6: Some of the visualized detection results obtained using our proposed RDAF-Net on the HRSC2016 dataset.

2) Results on DOTA: As shown in Table 5, our proposed method achieves the best detection performance compared to other advanced methods, reaching an mAP of 78.75\%. Through our proposed SI criterion, we divide the 15 target categories in the dataset into two typical categories: high aspect ratio targets (BR, GTF, LV, SH, and HA) and regularly shaped targets (PL, BD, SV, TC, BC, ST, SBF, RA, SP, and HC). The experimental results show that this method excels in the detection of high aspect ratio targets, achieving the best detection accuracy in all five subcategories; meanwhile, in the detection of regularly shaped targets, five out of ten subcategories have reached the best comparison level, which fully validates the robustness of our method in dealing with the diversity of target orientation and shape variability. It is worth noting that the improvement of detection accuracy for remote sensing targets with extreme aspect ratio features such as BR and LV is particularly significant, with an increase of 4.64\% and 4.59\% respectively compared to the second best method. The qualitative visualization results shown in Fig. 7 further demonstrate that this method has a significant advantage in detection effects under different scales, dense arrangement, and complex background conditions.

Figure 7: Some of the visualized detection results obtained using our proposed RDAF-Net on the DOTA dataset.

Table 5: Comparative experiments on the DOTA dataset.

Table 6: Comparative detection results on the UCAS-AOD dataset.

3) Results on UCAS-AOD: To further validate the effectiveness of the proposed RDAF-Net, a series of experiments were conducted on the UCAS-AOD dataset. The experimental results in Table 6 show that our method achieves the best performance among the compared detectors, with an mAP of 90.00\%. Specifically, the detection results for the categories of cars and airplanes are89.42\% and 90.57\%, respectively, both reaching the highest detection accuracy for these categories. This demonstrates that the proposed method exhibits strong robustness for densely arranged small objects, further validating its superior performance. Some qualitative results are shown in Figure 8.

Figure 8: Some of the visualization results of detection obtained using our proposed RDAF-Net on the UCAS-AOD dataset.

\section{ Conclusion }

This paper proposes a single-stage Rotated Dual Alignment Feature Network (RDAF-Net) for the problem of directional object detection in remote sensing images. The framework optimizes feature representation, anchor point refinement, and training sample selection, significantly improving detection performance. Specifically, RDAF-Net extracts high-quality orientation-sensitive features by adaptively adjusting convolution kernel parameters, effectively capturing the rotation characteristics of the target; secondly, by optimizing the anchor point generation mechanism, it ensures the spatial accuracy of the prediction box, significantly improving the feature alignment accuracy. Finally, in the sample selection strategy, combining target shape information and sample quality evaluation, it realizes the dynamic selection of positive samples, thereby enhancing the consistency between regression tasks and classification tasks. Experimental results show that RDAF-Net achieves excellent detection performance on three benchmark datasets: HRSC2016, DOTA, and UCAS-AOD, with mAP reaching 90.05\%, 79.60\%, and 90.00\% respectively, fully verifying the effectiveness and advancement of this method in the task of directional object detection in remote sensing images.

\section{ References }


\begin{figure}[t]
\centering
\includegraphics[width=0.45\textwidth]{fig_1.tiff}
\caption{Fig. 1}
\label{fig:fig_1}
\end{figure}

\begin{figure}[t]
\centering
\includegraphics[width=0.45\textwidth]{fig_2.tiff}
\caption{Fig. 2}
\label{fig:fig_2}
\end{figure}

\begin{figure}[t]
\centering
\includegraphics[width=0.45\textwidth]{fig_3.tiff}
\caption{Fig. 3}
\label{fig:fig_3}
\end{figure}

\begin{figure}[t]
\centering
\includegraphics[width=0.45\textwidth]{fig_4.tiff}
\caption{Fig. 4}
\label{fig:fig_4}
\end{figure}

\begin{figure}[t]
\centering
\includegraphics[width=0.45\textwidth]{fig_5.tiff}
\caption{Fig. 5}
\label{fig:fig_5}
\end{figure}

\begin{figure}[t]
\centering
\includegraphics[width=0.45\textwidth]{fig_6.tiff}
\caption{Fig. 6}
\label{fig:fig_6}
\end{figure}

\begin{figure}[t]
\centering
\includegraphics[width=0.45\textwidth]{fig_7.tiff}
\caption{Fig. 7}
\label{fig:fig_7}
\end{figure}

\begin{figure}[t]
\centering
\includegraphics[width=0.45\textwidth]{fig_8.tif}
\caption{Fig. 8}
\label{fig:fig_8}
\end{figure}

\begin{table}[t]
\centering
\begin{tabular}{c c c c c}
\toprule
\textbf{} & \textbf{Baseline} & \textbf{Different Variants} & \textbf{Different Variants} & \textbf{Different Variants} \\
\midrule
DRRCM &  & √ & √ & √ \\
ARFAM &  & × & √ & √ \\
SAQA &  & × & × & √ \\
mAP & 71.64 & 73.43 & 84.89 & 89.58 \\
\bottomrule
\end{tabular}
\caption{Table 1}
\label{tab:tbl_1}
\end{table}

\begin{table}[t]
\centering
\begin{tabular}{c c c}
\toprule
\textbf{Method} & \textbf{Backbone} & \textbf{mAP} \\
\midrule
RetinaNet[12] & R50 & 71.64 \\
RetinaNet[12] & R-50- DRRCM\_S4 & 71.93 \\
RetinaNet[12] & R-50- DRRCM\_S4+S3 & 72.66 \\
RetinaNet[12] & R-50- DRRCM\_S4+S3+S2 & 73.43 \\
\bottomrule
\end{tabular}
\caption{Table 2}
\label{tab:tbl_2}
\end{table}

\begin{table}[t]
\centering
\begin{tabular}{c c c c c}
\toprule
\textbf{Method} & \textbf{DRRCM} & \textbf{ARFAM} & \textbf{SAQA} & \textbf{mAP} \\
\midrule
RDAF-Net & R-50- DRRCM\_S4+S3+S2 & 1 & √ & 88.74 \\
RDAF-Net & R-50- DRRCM\_S4+S3+S2 & 2 & √ & 89.58 \\
RDAF-Net & R-50- DRRCM\_S4+S3+S2 & 3 & √ & 89.43 \\
RDAF-Net & R-50- DRRCM\_S4+S3+S2 & 4 & √ & 88.94 \\
\bottomrule
\end{tabular}
\caption{Table 3}
\label{tab:tbl_3}
\end{table}

\begin{table}[t]
\centering
\begin{tabular}{c c c c c}
\toprule
\textbf{Methods} & \textbf{Backbone} & \textbf{Size} & \textbf{Anchor} & \textbf{mAP} \\
\midrule
R2CNN[41] & R-101 & 800×800 & 21 & 73.07 \\
RRPN[21] & R-101 & 800×800 & 54 & 79.08 \\
RRD[42] & VGG16 & 384×384 & 13 & 84.30 \\
RoI-Trans.[13] & R-101 & 512×800 & 5 & 86.20 \\
DAL[30] & R-101 & 416×416 & 3 & 88.95 \\
R-RetinaNet[12] & R-101 & 800×800 & 121 & 89.18 \\
R3Det[11] & R-101 & 800×800 & 126 & 89.26 \\
CFC-Net[43] & R-101 & 800×800 & 1 & 89.50 \\
RDAF-Net & R-50-DRRCM & 512×800 & 1 & 89.58 \\
RDAF-Net（aug） & R-50-DRRCM & 512×800 & 1 & 90.03 \\
RDAF-Net（aug） & R-101-DRRCM & 512×800 & 1 & 90.05 \\
\bottomrule
\end{tabular}
\caption{Table 4}
\label{tab:tbl_4}
\end{table}

\begin{table}[t]
\centering
\begin{tabular}{c c c c c c c c c c c c c c c c c c}
\toprule
\textbf{Methods} & \textbf{Backbone} & \textbf{PL} & \textbf{BD} & \textbf{BR} & \textbf{GTF} & \textbf{SV} & \textbf{LV} & \textbf{SH} & \textbf{TC} & \textbf{BC} & \textbf{ST} & \textbf{SBF} & \textbf{RA} & \textbf{HA} & \textbf{SP} & \textbf{HC} & \textbf{mAP} \\
\midrule
ICN[22] & R-101 & 81.36 & 74.30 & 47.70 & 70.32 & 64.89 & 67.82 & 69.98 & 90.76 & 79.06 & 78.20 & 53.64 & 62.90 & 67.02 & 64.17 & 50.23 & 68.16 \\
RoI-Trans.[13] & R-101 & 88.64 & 78.52 & 43.44 & 75.92 & 68.81 & 73.68 & 83.59 & 90.74 & 77.27 & 81.46 & 58.39 & 53.54 & 62.83 & 58.93 & 47.67 & 69.56 \\
CAD-Net[44] & R-101 & 87.80 & 82.40 & 49.40 & 73.50 & 71.10 & 63.50 & 76.70 & 90.90 & 79.20 & 73.30 & 48.40 & 60.90 & 62.00 & 67.00 & 62.20 & 69.90 \\
DRN[45] & H-104 & 88.91 & 80.22 & 43.52 & 63.35 & 73.48 & 70.69 & 84.94 & 90.14 & 83.85 & 84.11 & 50.12 & 58.41 & 67.62 & 68.60 & 52.50 & 70.70 \\
O2-DNet[46] & H-104 & 89.31 & 82.14 & 47.33 & 61.21 & 71.32 & 74.03 & 78.62 & 90.76 & 82.23 & 81.36 & 60.93 & 60.17 & 58.21 & 66.98 & 61.03 & 71.04 \\
R3Det[11] & R-101 & 89.54 & 81.99 & 48.46 & 62.52 & 70.48 & 74.29 & 77.54 & 90.80 & 81.39 & 83.54 & 61.97 & 59.82 & 65.44 & 67.46 & 60.05 & 71.69 \\
SCRDet[47] & R-101 & 89.98 & 80.65 & 52.09 & 68.36 & 68.36 & 60.32 & 72.41 & 90.85 & 87.94 & 86.86 & 65.02 & 66.68 & 66.25 & 68.24 & 65.21 & 72.61 \\
FADet[48] & R-101 & 90.21 & 79.58 & 45.49 & 76.41 & 73.18 & 68.27 & 79.56 & 90.83 & 83.40 & 84.64 & 53.40 & 65.42 & 74.17 & 69.69 & 64.86 & 73.28 \\
CSL-F[17] & R-152 & 90.25 & 85.53 & 54.64 & 75.31 & 70.44 & 73.51 & 77.62 & 90.84 & 86.15 & 86.69 & 69.60 & 68.04 & 73.83 & 71.10 & 68.93 & 76.17 \\
DAL[30] & R-50 & 89.69 & 83.11 & 55.03 & 71.00 & 78.30 & 81.90 & 88.46 & 90.89 & 84.97 & 87.46 & 64.41 & 65.65 & 76.86 & 72.09 & 64.35 & 76.95 \\
RDAF-Net
(aug) & R-101-DRRCM & 89.99 & 75.77 & 49.94 & 77.83 & 70.24 & 82.41 & 89.08 & 90.71 & 66.62 & 88.46 & 73.07 & 65.04 & 71.97 & 57.95 & 42.36 & 72.76 \\
RDAF-Net
(aug+ms) & R-101-DRRCM & 90.27 & 85.24 & 59.60 & 85.52 & 78.66 & 87.02 & 89.96 & 90.68 & 75.18 & 89.75 & 76.93 & 74.63 & 77.66 & 70.59 & 49.65 & 78.75 \\
\bottomrule
\end{tabular}
\caption{Table 5}
\label{tab:tbl_5}
\end{table}

\begin{table}[t]
\centering
\begin{tabular}{c c c c}
\toprule
\textbf{Methods} & \textbf{Car} & \textbf{Airplane} & \textbf{mAP} \\
\midrule
YOLOv3[49] & 74.63 & 89.52 & 82.08 \\
R-RetinaNet[12] & 84.64 & 90.51 & 87.57 \\
Faster RCNN[50] & 86.87 & 89.86 & 88.36 \\
RoI-Trans.[13] & 88.02 & 90.02 & 89.02 \\
RIDet-O[51] & 88.88 & 90.35 & 89.62 \\
RDAF-Net & 89.42 & 90.57 & 90.00 \\
\bottomrule
\end{tabular}
\caption{Table 6}
\label{tab:tbl_6}
\end{table}

\begin{table}[t]
\centering
\begin{tabular}{c c}
\toprule
\textbf{Airplane



Car} & \textbf{} \\
\midrule
\bottomrule
\end{tabular}
\caption{Table 7}
\label{tab:tbl_7}
\end{table}

\begin{equation}
SI = \frac{\min(w, h)}{\max(w, h)}
\label{eq:eq_1}
\end{equation}

\begin{equation}
\frac{1}{2}
\label{eq:eq_2}
\end{equation}

\begin{equation}
\frac{1}{2}
\label{eq:eq_3}
\end{equation}

\begin{equation}
P_{\text{avg}}(\cdot)
\label{eq:eq_4}
\end{equation}

\begin{equation}
P_{\text{max}}(\cdot)
\label{eq:eq_5}
\end{equation}

\begin{equation}
S_{ag} = P_{ag}(A), S_{max} = P_{max}(A)
\label{eq:eq_6}
\end{equation}

\begin{equation}
F^2 \to C_{\text{in}} (\cdot)
\label{eq:eq_7}
\end{equation}

\begin{equation}
P_{\text{avg}}(\cdot)
\label{eq:eq_8}
\end{equation}

\begin{equation}
P_{\text{max}}(\cdot)
\label{eq:eq_9}
\end{equation}

\begin{equation}
S_{ag} = P_{ag}(A), S_{max} = P_{max}(A)
\label{eq:eq_10}
\end{equation}

\begin{equation}
F^2 \to C_{\text{in}} (\cdot)
\label{eq:eq_11}
\end{equation}

\begin{equation}
S = F^2 C_n \left( \left[ S_{\text{min}} ; S_{\text{max}} \right] \right)
\label{eq:eq_12}
\end{equation}

\begin{equation}
S'
\label{eq:eq_13}
\end{equation}

\begin{equation}
S_{i'}
\label{eq:eq_14}
\end{equation}

\begin{equation}
S_i' = \text{sigmoid}(S')
\label{eq:eq_15}
\end{equation}

\begin{equation}
V c_{in}
\label{eq:eq_16}
\end{equation}

\begin{equation}
V_{c,\dot{m}} = P_{\text{avg}} (S_i \times F)
\label{eq:eq_17}
\end{equation}

\begin{equation}
\theta_i
\label{eq:eq_18}
\end{equation}

\begin{equation}
\theta_i = K(\text{Softsign}(z_\theta))
\label{eq:eq_19}
\end{equation}

\begin{equation}
Z \theta
\label{eq:eq_20}
\end{equation}

\begin{equation}
K = \left( \frac{p}{180.0} \times \pi \right)
\label{eq:eq_21}
\end{equation}

\begin{equation}
p
\label{eq:eq_22}
\end{equation}

\begin{equation}
\lambda_i
\label{eq:eq_23}
\end{equation}

\begin{equation}
\lambda_i = \text{sigmoid}(z_a)
\label{eq:eq_24}
\end{equation}

\begin{equation}
Z_a
\label{eq:eq_25}
\end{equation}

\begin{equation}
\theta_i
\label{eq:eq_26}
\end{equation}

\begin{equation}
Y'_i = rotate(Y_i, -\theta_i)
\label{eq:eq_27}
\end{equation}

\begin{equation}
W_1-interpolatio(W_1, Y_1)
\label{eq:eq_28}
\end{equation}

\begin{equation}
Y_i
\label{eq:eq_29}
\end{equation}

\begin{equation}
Y_i'
\label{eq:eq_30}
\end{equation}

\begin{equation}
Y_i
\label{eq:eq_31}
\end{equation}

\begin{equation}
\theta_i
\label{eq:eq_32}
\end{equation}

\begin{equation}
W_i'
\label{eq:eq_33}
\end{equation}

\begin{equation}
\text{interpolation}(\cdot)
\label{eq:eq_34}
\end{equation}

\begin{equation}
W_i
\label{eq:eq_35}
\end{equation}

\begin{equation}
Y = F \cdot \sum_{i=1}^n \lambda_i W_i'
\label{eq:eq_36}
\end{equation}

\begin{equation}
\Delta \theta_g = \frac{\theta_g - \theta}{\pi} + k
\label{eq:eq_37}
\end{equation}

\begin{equation}
\Delta x_g = \frac{(x_g - x) \cdot R(\theta)}{w}
\label{eq:eq_38}
\end{equation}

\begin{equation}
\Delta y_g = \frac{(y_g - x) \cdot R(\theta)}{h}
\label{eq:eq_39}
\end{equation}

\begin{equation}
\Delta w_g = \log\left( \frac{w_g}{w} \right)
\label{eq:eq_40}
\end{equation}

\begin{equation}
\Delta h_g = \log\left( \frac{h^g}{h} \right)
\label{eq:eq_41}
\end{equation}

\begin{equation}
\Delta \theta_g = \frac{\theta_g - \theta}{\pi} + k
\label{eq:eq_42}
\end{equation}

\begin{equation}
0
\label{eq:eq_43}
\end{equation}

\begin{equation}
0 = \sum_{p_n \in R} \left( L_p^n - (p_0 + p_n) \right)
\label{eq:eq_44}
\end{equation}

\begin{equation}
L^n_p
\label{eq:eq_45}
\end{equation}

\begin{equation}
(p_0 + p_n)
\label{eq:eq_46}
\end{equation}

\begin{equation}
p^\circ
\label{eq:eq_47}
\end{equation}

\begin{equation}
p_n
\label{eq:eq_48}
\end{equation}

\begin{equation}
\left[ (p_x, p_y) \right]
\label{eq:eq_49}
\end{equation}

\begin{equation}
\mathrm{RK}(1)(10)_(0)(0)
\label{eq:eq_50}
\end{equation}

\begin{equation}
Y(p_0) = \sum_{p_n \in R} w(p_n) \cdot x(p_0 + p_n + \theta)
\label{eq:eq_51}
\end{equation}

\begin{equation}
Y(p)
\label{eq:eq_52}
\end{equation}

\begin{equation}
w(p_n)
\label{eq:eq_53}
\end{equation}

\begin{equation}
p_n
\label{eq:eq_54}
\end{equation}

\begin{equation}
x(.)
\label{eq:eq_55}
\end{equation}

\begin{equation}
D_i^n = (\mu + \sigma) e^{-1/(2(g - a))}
\label{eq:eq_56}
\end{equation}

\begin{equation}
D_i^{\text{IoU}}
\label{eq:eq_57}
\end{equation}

\begin{equation}
\mu
\label{eq:eq_58}
\end{equation}

\begin{thebibliography}{99}

\bibitem{ref_1}
\textit{1.	Cheng, Gong, Peicheng Zhou, and Junwei Han. "Learning Rotation-Invariant Convolutional Neural Net}.

\bibitem{ref_2}
\textit{2.	Zaidi, Syed Sahil Abbas, Mohammad Samar Ansari, Asra Aslam, Nadia Kanwal, Mamoona Asghar, and Bri}.

\bibitem{ref_3}
\textit{3.	Zou, Zhengxia, Keyan Chen, Zhenwei Shi, Yuhong Guo, and Jieping Ye. "Object Detection in 20 Years}.

\bibitem{ref_4}
\textit{4.	Zheng, Zhe, Lin Lei, Hao Sun, and Gangyao Kuang. "A Review of Remote Sensing Image Object Detecti}.

\bibitem{ref_5}
\textit{5.	Yao, Qunli, Xian Hu, and Hong Lei. "Geospatial Object Detection in Remote Sensing Images Based on}.

\bibitem{ref_6}
\textit{6.	Bai, Peng, Ying Xia, and Jiangfan Feng. "Composite Perception and Multi-Scale Fusion Network for }.

\bibitem{ref_7}
\textit{7.	Deng, Zhipeng, Hao Sun, Shilin Zhou, Juanping Zhao, Lin Lei, and Huanxin Zou. "Multi-Scale Object}.

\bibitem{ref_8}
\textit{8.	Agarwal, Shivang, Jean Ogier Du Terrail, and Frédéric Jurie. "Recent Advances in Object Detection}.

\bibitem{ref_9}
\textit{9.	Wang, Guanqun, Yin Zhuang, He Chen, Xiang Liu, Tong Zhang, Lianlin Li, Shan Dong, and Qianbo Sang}.

\bibitem{ref_10}
\textit{10.	Zhou, Yang, and Hui Li. "A Survey of Dense Object Detection Methods Based on Deep Learning." IEE}.

\bibitem{ref_11}
\textit{11.	Yang, Xue, Junchi Yan, Ziming Feng, and Tao He. "R3det: Refined Single-Stage Detector with Featu}.

\bibitem{ref_12}
\textit{12.	Lin, Tsung-Yi, Priya Goyal, Ross Girshick, Kaiming He, and Piotr Dollar. "Focal Loss for Dense O}.

\bibitem{ref_13}
\textit{13.	Ding, Jian, Nan Xue, Yang Long, Gui-Song Xia, and Qikai Lu. "Learning Roi Transformer for Orient}.

\bibitem{ref_14}
\textit{14.	Xu, Yongchao, Mingtao Fu, Qimeng Wang, Yukang Wang, Kai Chen, Gui-Song Xia, and Xiang Bai. "Glid}.

\bibitem{ref_15}
\textit{15.	Cheng, Yuhu, Chengqing Xu, Yi Kong, and Xuesong Wang. "Short-Side Excursion for Oriented Object }.

\bibitem{ref_16}
\textit{16.	Xie, Xingxing, Gong Cheng, Jiabao Wang, Xiwen Yao, and Junwei Han. "Oriented R-Cnn for Object De}.

\bibitem{ref_17}
\textit{17.	Yang, Xue, and Junchi Yan. "Arbitrary-Oriented Object Detection with Circular Smooth Label." Pap}.

\bibitem{ref_18}
\textit{18.	Yang, Xue, Xiaojiang Yang, Jirui Yang, Qi Ming, Wentao Wang, Qi Tian, and Junchi Yan. "Learning }.

\bibitem{ref_19}
\textit{19.	Yang, Xue, Yue Zhou, Gefan Zhang, Jirui Yang, Wentao Wang, Junchi Yan, Xiaopeng Zhang, and Qi Ti}.

\bibitem{ref_20}
\textit{20.	Shen, Chongchong, Jiangbo Qian, Chong Wang, Diqun Yan, and Caiming Zhong. "Dynamic Sensing and C}.

\bibitem{ref_21}
\textit{21.	Ma, Jianqi, Weiyuan Shao, Hao Ye, Li Wang, Hong Wang, Yingbin Zheng, and Xiangyang Xue. "Arbitra}.

\bibitem{ref_22}
\textit{22.	Azimi, Seyed Majid, Eleonora Vig, Reza Bahmanyar, Marco Körner, and Peter Reinartz. "Towards Mul}.

\bibitem{ref_23}
\textit{23.	Liu, Xulun, Shiping Ma, Linyuan He, Chen Wang, and Zhe Chen. "Hybrid Network Model: Transconvnet}.

\bibitem{ref_24}
\textit{24.	Liu, Ze, Yutong Lin, Yue Cao, Han Hu, Yixuan Wei, Zheng Zhang, Stephen Lin, and Baining Guo. "Sw}.

\bibitem{ref_25}
\textit{25.	Sun, Peng, Yongbin Zheng, Zongtan Zhou, Wanying Xu, and Qiang Ren. "R4 Det: Refined Single-Stage}.

\bibitem{ref_26}
\textit{26.	Zheng, Yongbin, Peng Sun, Zongtan Zhou, Wanying Xu, and Qiang Ren. "Adt-Det: Adaptive Dynamic Re}.

\bibitem{ref_27}
\textit{27.	Zhang, Shifeng, Cheng Chi, Yongqiang Yao, Zhen Lei, and Stan Z Li. "Bridging the Gap between Anc}.

\bibitem{ref_28}
\textit{28.	Ma, Wenping, Xiaoteng Wang, Hao Zhu, Xiaoting Yang, Xiaoyu Yi, and Licheng Jiao. "Significant Fe}.

\bibitem{ref_29}
\textit{29.	Li, Wentong, Yijie Chen, Kaixuan Hu, and Jianke Zhu. "Oriented Reppoints for Aerial Object Detec}.

\bibitem{ref_30}
\textit{30.	Ming, Qi, Zhiqiang Zhou, Lingjuan Miao, Hongwei Zhang, and Linhao Li. "Dynamic Anchor Learning f}.

\bibitem{ref_31}
\textit{31.	Zhang, Xiaosong, Fang Wan, Chang Liu, Rongrong Ji, and Qixiang Ye. "Freeanchor: Learning to Matc}.

\bibitem{ref_32}
\textit{32.	He, Yihui, Chenchen Zhu, Jianren Wang, Marios Savvides, and Xiangyu Zhang. "Bounding Box Regress}.

\bibitem{ref_33}
\textit{33.	Jiang, Borui, Ruixuan Luo, Jiayuan Mao, Tete Xiao, and Yuning Jiang. "Acquisition of Localizatio}.

\bibitem{ref_34}
\textit{34.	Feng, Di, Lars Rosenbaum, and Klaus Dietmayer. "Towards Safe Autonomous Driving: Capture Uncerta}.

\bibitem{ref_35}
\textit{35.	Liu, Zikun, Liu Yuan, Lubin Weng, and Yiping Yang. "A High Resolution Optical Satellite Image Da}.

\bibitem{ref_36}
\textit{36.	Xia, Gui-Song, Xiang Bai, Jian Ding, Zhen Zhu, Serge Belongie, Jiebo Luo, Mihai Datcu, Marcello }.

\bibitem{ref_37}
\textit{37.	Zhu, Haigang, Xiaogang Chen, Weiqun Dai, Kun Fu, Qixiang Ye, and Jianbin Jiao. "Orientation Robu}.

\bibitem{ref_38}
\textit{38.	He, Kaiming, Xiangyu Zhang, Shaoqing Ren, and Jian Sun. "Deep Residual Learning for Image Recogn}.

\bibitem{ref_39}
\textit{39.	Everingham, Mark, Luc Van Gool, Christopher KI Williams, John Winn, and Andrew Zisserman. "The P}.

\bibitem{ref_40}
\textit{40.	Chen, Kai, Jiaqi Wang, Jiangmiao Pang, Yuhang Cao, Yu Xiong, Xiaoxiao Li, Shuyang Sun, Wansen Fe}.

\bibitem{ref_41}
\textit{41.	Jiang, Yingying, Xiangyu Zhu, Xiaobing Wang, Shuli Yang, Wei Li, Hua Wang, Pei Fu, and Zhenbo Lu}.

\bibitem{ref_42}
\textit{42.	Liao, Minghui, Zhen Zhu, Baoguang Shi, Gui-song Xia, and Xiang Bai. "Rotation-Sensitive Regressi}.

\bibitem{ref_43}
\textit{43.	Ming, Qi, Lingjuan Miao, Zhiqiang Zhou, and Yunpeng Dong. "Cfc-Net: A Critical Feature Capturing}.

\bibitem{ref_44}
\textit{44.	Zhang, Gongjie, Shijian Lu, and Wei Zhang. "Cad-Net: A Context-Aware Detection Network for Objec}.

\bibitem{ref_45}
\textit{45.	Pan, Xingjia, Yuqiang Ren, Kekai Sheng, Weiming Dong, Haolei Yuan, Xiaowei Guo, Chongyang Ma, an}.

\bibitem{ref_46}
\textit{46.	Wei, Haoran, Yue Zhang, Zhonghan Chang, Hao Li, Hongqi Wang, and Xian Sun. "Oriented Objects as }.

\bibitem{ref_47}
\textit{47.	Yang, Xue, Jirui Yang, Junchi Yan, Yue Zhang, Tengfei Zhang, Zhi Guo, Xian Sun, and Kun Fu. "Scr}.

\bibitem{ref_48}
\textit{48.	Li, Chengzheng, Chunyan Xu, Zhen Cui, Dan Wang, Tong Zhang, and Jian Yang. "Feature-Attentioned }.

\bibitem{ref_49}
\textit{49.	Farhadi, Ali, and Joseph Redmon. "Yolov3: An Incremental Improvement." Paper presented at the Co}.

\bibitem{ref_50}
\textit{50.	Ren, Shaoqing, Kaiming He, Ross Girshick, and Jian Sun. "Faster R-Cnn: Towards Real-Time Object }.

\bibitem{ref_51}
\textit{51.	Ming, Qi, Lingjuan Miao, Zhiqiang Zhou, Xue Yang, and Yunpeng Dong. "Optimization for Arbitrary-}.

\end{thebibliography}
\end{document}
